
\subsubsection{2.1 Parameter-Efficient Fine-Tuning for SSMs}

SSM-PEFT provides analysis of fine-tuning methods for State Space Models, introducing Sparse Dimension Tuning (SDT) as an alternative to standard LoRA. MambaPEFT catalogs over 20 PEFT method variants applied to Mamba architectures, finding that PEFT is generally more effective for Mamba than for Transformers on certain tasks. State-offset Tuning proposes learning an additive offset to the SSM state. These works identify successful and unsuccessful configurations but do not provide a measurable criterion for predicting when projection-only adaptation will succeed.

\subsubsection{2.2 State Space Model Foundations}

Mamba introduced selective state spaces with input-dependent discretization, achieving Transformer-competitive performance with linear-time inference. The architecture parameterizes the continuous-time state matrix A through A\_log, with discretization via the zero-order hold: A\_bar = exp(Delta * A). While the discretization step Delta is input-dependent, the A matrix itself is fixed after pretraining.

\subsubsection{2.3 Eigenvalue Analysis in Recurrent Models}

The connection between eigenvalue magnitude and memory capacity in recurrent systems is established in control theory. For discrete-time linear systems h\_t = A \textit{ $h_{t-1}$ + B } u\_t, eigenvalues |lambda| < 1 ensure stability, with |lambda| approaching 1 corresponding to longer memory. Prior work has not applied eigenvalue analysis to understand PEFT limitations in SSMs.
